% TOP_PROBLEM: {"id": 30007046, "title": "Rational effectivity modulo a controlled boundary", "category_id": 5, "category": {"id": 5, "name": "algebraic_geometry", "display_name": "Algebraic Geometry", "description": "Geometric objects defined by polynomial equations.", "slug": "algebraic-geometry", "order_index": 5, "created_at": "2026-07-31T15:26:25.671Z"}, "status": "open"}
% FIRST_POSTED: 2026-10-06
% ATTEMPT_STATUS: solved
% OPENAI_CLAIM: {"schema_version":1,"source_repo":"https://github.com/openai/math","source_commit":"adc7f1241b42e322a6451854ab7e4b4c146bf78a","release_date":"2026-10-06","match":"full","resolution":"proved","independently_reviewed":false,"second_pass":"agreed","families":[{"family":"068","title":"Anticanonical nonvanishing in every dimension","subject":"Algebraic and complex geometry","manuscripts":[{"title":"Integrable metrics and effectivity with controlled boundary","pdf_path":"preprints/Integrable-metrics-and-effectivity-with-controlled-boundary-September-26-2026/main.pdf","readme_path":"preprints/Integrable-metrics-and-effectivity-with-controlled-boundary-September-26-2026/README.md","date":"2026-09-26","theorem_ref":"thm:controlled-boundary (source/sections/introduction.tex:24); conv:fixed-error (source/sections/conversion.tex:10)","statement_sources":[{"path":"preprints/Integrable-metrics-and-effectivity-with-controlled-boundary-September-26-2026/build/source/sections/introduction.tex","line":24,"label":"thm:controlled-boundary"},{"path":"preprints/Integrable-metrics-and-effectivity-with-controlled-boundary-September-26-2026/build/source/sections/conversion.tex","line":10,"label":"conv:fixed-error"}]}],"lean":null,"reasoning_trace":null}]}
\documentclass[11pt]{article}
\usepackage[margin=25mm]{geometry}
\usepackage{fontspec}
\setmainfont{Latin Modern Roman}
\usepackage{amsmath,amssymb,mathtools}
\usepackage{unicode-math}
\setmathfont{Latin Modern Math}
\usepackage{enumitem}
\usepackage{xurl,hyperref,bookmark}
\hypersetup{colorlinks=true,urlcolor=blue}
\setcounter{secnumdepth}{0}
\setlength{\parindent}{0pt}
\setlength{\parskip}{5pt}
\setlength{\emergencystretch}{3em}
\begin{document}
\section*{Rational effectivity modulo a controlled boundary}

\subsection{OpenAI's claim}
A globally strict anticanonical metric with controlled boundary integrability is claimed to make every pseudo-effective rational divisor effective after a correction on that boundary, and to remove fixed errors from unbounded anticanonical twists. The cited manuscript argues for that result; its relation to this problem is specified below.

\subsection{Mathematical statement of the claimed result}
\paragraph{Mathematical claim 1.} Let $Y$ be a complex projective variety that is smooth and has positive dimension and $C\ge0$ integral. Assume $-K_Y+C$ carries a singular Hermitian metric $h$ with curvature at least a Kähler form and whose multiplier ideal satisfies $\mathcal J(h)\supset\mathcal O_Y(-C)$. OpenAI claims
\[
\forall\text{ pseudo-effective rational }M\ \exists C_0\ge0\text{ supported on }C:\quad M+C_0\sim_{\mathbb Q}E\ge0.
\]
The multiplier condition means local squared integrability of the canonical section of $C$ in the metric. A separate conversion theorem says that if $X$ is smooth connected projective complex, $L=-K_X$ smoothly semipositive and $P$ pseudo-effective Cartier, then effective integral $N_j\sim m_jL-P$ for unbounded positive $m_j$ imply $H^0(X,mL)\ne0$ for some $m>0$.

\subsection{Outline of the argument}
Manuscript 1 (family 068, 2026-09-26): Singular adjoint direct-image positivity descends a strict metric with an exceptional integral boundary and trivial total-space multiplier ideal. Controlled section rings and a single-Iitaka conversion return the corrected sections upstairs; a supported Fano-type base model supplies the required effectivity without losing the disappearing-prime condition.

\subsection{Scope relative to this problem}
Manuscript 1 (family 068, 2026-09-26): The problem selects the mathematical statement summarized here, including its stated hypotheses and parameter range. The cited theorem gives a claimed proof of that assertion in full.

\subsection{Formal verification}
Family 068 has no Lean scope page in this snapshot; no formalization of these manuscript results is listed.

This site has not executed the Lean files or checked their dependencies. These descriptions report the scope of the linked declarations, rather than a new verification of the proofs.

\subsection{Status caveat}
OpenAI attributes these results to an internal model. Its release notice cautions that results lacking full formalization can contain mistakes. This site has not independently assessed the proofs; a solved label records the scope of the reported claim.

\subsection{Sources}
\begin{itemize}
\item OpenAI, \emph{Integrable metrics and effectivity with controlled boundary}, 2026-09-26; citation key \texttt{OAI:Integrable-metrics-and-effectivity-with-controlled-boundary-September-26-2026}. \href{https://github.com/openai/math/blob/adc7f1241b42e322a6451854ab7e4b4c146bf78a/preprints/Integrable-metrics-and-effectivity-with-controlled-boundary-September-26-2026/main.pdf}{PDF (GitHub)}; \href{https://github.com/openai/math/raw/adc7f1241b42e322a6451854ab7e4b4c146bf78a/preprints/Integrable-metrics-and-effectivity-with-controlled-boundary-September-26-2026/main.pdf}{PDF (direct)}; \href{https://github.com/openai/math/blob/adc7f1241b42e322a6451854ab7e4b4c146bf78a/preprints/Integrable-metrics-and-effectivity-with-controlled-boundary-September-26-2026/README.md}{Manuscript page and citation}; \href{https://github.com/openai/math/blob/main/preprints/Integrable-metrics-and-effectivity-with-controlled-boundary-September-26-2026/main.pdf}{Latest version}; Statement source: \href{https://github.com/openai/math/blob/adc7f1241b42e322a6451854ab7e4b4c146bf78a/preprints/Integrable-metrics-and-effectivity-with-controlled-boundary-September-26-2026/build/source/sections/introduction.tex\#L24}{thm:controlled-boundary}; Statement source: \href{https://github.com/openai/math/blob/adc7f1241b42e322a6451854ab7e4b4c146bf78a/preprints/Integrable-metrics-and-effectivity-with-controlled-boundary-September-26-2026/build/source/sections/conversion.tex\#L10}{conv:fixed-error}.
\item \href{https://github.com/openai/math/blob/adc7f1241b42e322a6451854ab7e4b4c146bf78a/CONTENTS.md}{OpenAI catalogue entry}.
\end{itemize}
\end{document}
